\PassOptionsToPackage{table,xcdraw}{xcolor}
\documentclass[11pt, a4paper]{article}

\usepackage[T1]{fontenc}
\usepackage[utf8]{inputenc}
\usepackage{tgpagella}
\usepackage[spanish, es-noshorthands]{babel}
\usepackage{amsmath, amssymb}
\usepackage[most]{tcolorbox}
\usepackage{tabularx}
\usepackage[american]{circuitikz}
\usepackage[margin=2.5cm]{geometry}
\usepackage{fancyhdr}

\pagestyle{fancy}
\fancyhf{}
\setlength{\headheight}{16pt}
\lhead{\textbf{UCB Sede Tarija} | Ing. Mecatrónica}
\rhead{IMT-221 | Laboratorio Diferenciado S08-2}
\renewcommand{\headrulewidth}{0.4pt}
\definecolor{ucbblue}{RGB}{0, 51, 102}

\begin{document}

\begin{tcolorbox}[colback=ucbblue!5, colframe=ucbblue, title=\textbf{Guía de Laboratorio: Checkpoints Diferenciados (Hito 2)}]
    \textbf{Regla de Operación:} El trabajo continúa desde el último checkpoint validado de su grupo. Determine su estado actual y proceda estrictamente según el flujo indicado.
\end{tcolorbox}

\section*{Evaluación de Estado Actual (Triage)}
Marque la casilla correspondiente al estado validado por el docente en la sesión anterior:
\begin{itemize}
    \item[$\square$] \textbf{Estado A (Espejo Pendiente):} No hay validación de la fuente de cola.
    \item[$\square$] \textbf{Estado B (Espejo Validado / Par Pendiente):} $I_T$ demostrada, pero Par Q1-Q2 no medido.
    \item[$\square$] \textbf{Estado C (Par Validado / Diferencial Pendiente):} Punto Q aprobado, falta inyectar señal.
    \item[$\square$] \textbf{Estado D (Respuesta Diferencial Disponible):} Preparado para Modo Común y CMRR.
\end{itemize}

\vspace{0.3cm}
\hrule
\vspace{0.3cm}

\section*{Rutas de Acción según su Estado}

\subsection*{Para Estado A: Espejo de Corriente}
\textbf{Objetivo:} Validar la fuente de cola antes de conectar Q1 y Q2.
\begin{enumerate}
    \item Verifique magnitud y polaridad de los rieles ($+9\,\text{V}, 0\,\text{V}, -9\,\text{V}$) y pinout en el datasheet.
    \item Mida $I_{REF}$ a través de $R_{REF}$ y verifique $V_{BE}$ de la referencia diódica Q3[cite: 5].
    \item Utilice una carga temporal ($R_{TEST}$) en Q4 para inferir $I_T$. Solicite el \textbf{GO} del docente[cite: 5].
\end{enumerate}

\subsection*{Para Estado B: Par Diferencial Balanceado}
\textbf{Objetivo:} Encontrar el Punto de Operación estático con $v_1 = v_2 = 0\,\text{V}$[cite: 5].
\begin{enumerate}
    \item Conecte Q1 y Q2 con las cargas $R_{C1}$ y $R_{C2}$ y emisores unidos al colector de Q4[cite: 5].
    \item Mida $v_E$, $v_{C1}$, $v_{C2}$ y verifique la operación en \textbf{Activa Directa} ($V_{CE} > 0.3\,\text{V}$, $V_{BC} < 0$)[cite: 5].
    \item Calcule el Offset Estático: $V_{OD,0} = v_{C2} - v_{C1}$. Solicite el \textbf{GO} del docente[cite: 5].
\end{enumerate}

\subsection*{Para Estado C: Ganancia Diferencial ($A_d$) y Shunt}
\textbf{Objetivo:} Predecir e inyectar señal diferencial pura[cite: 5].
\begin{enumerate}
    \item Calcule $g_m = I_{CQ,\text{medido}} / V_T$. Prediga $A_d = g_m R_C$ y la salida esperada $v_{od, \text{pred}}$[cite: 5].
    \item Arme la rama del Shunt ($R_{LIM}, R_{sh}$) y mida el voltaje generado $v_d = v_1 - v_2$ real[cite: 5].
    \item Inyecte $v_d$ al par diferencial y mida la salida $v_{od} = v_{C2} - v_{C1}$ con osciloscopio o DMM de precisión. Compare con su predicción teórica[cite: 5].
\end{enumerate}

\subsection*{Para Estado D: Preparación de Modo Común ($v_{cm}$)}
\textbf{Objetivo:} Diseñar la medición de la ganancia común ($A_c$)[cite: 5].
\begin{enumerate}
    \item \textbf{Defina su Convención:} Declare explícitamente si usará salida Single-Ended o Diferencial. \textit{Advertencia: CMRR exige usar la misma convención para $A_d$ y $A_c$}[cite: 5].
    \item Realice la simulación de Modo Común ($v_1 = v_2 = v_{cm}$) en LTspice y genere su plan de medición antes de inyectar señal física[cite: 5].
\end{enumerate}

\end{document}
